\chapter[A priori bounds and limbs]{Source digests: a priori bounds, satellites, Feigenbaum points, and limbs}
\label{ch:digests-apriori-limbs}

This chapter digests the sources that give the strongest current results on local
connectivity of $\M$ beyond the foundations of Chapters~\ref{ch:digests-foundations}
and~\ref{ch:digests-modern}: the Kahn--Lyubich series on a priori bounds (Parts~I--IV),
the Dudko--Lyubich papers on satellite parameters and on Feigenbaum points, the
Kahn--Kapiamba--Lyubich paper on parabolically bounded primitive renormalization, and
the Dudko--Schleicher paper on limbs.  The conventions are those of the earlier digest
chapters.  Definitions and theorems are the source's own, numbered as the source prints
them.  When a source does not print its main statement as a numbered theorem, we say so.
Each digest ends with its MLC scope, which is the information used in
Chapter~\ref{ch:approaches-map}.

\section{Common framework}
\label{sec:apriori-framework}

A \emph{quadratic-like map} is a proper holomorphic map $g\colon U\to V$ of degree two
between Jordan disks with $\overline U\subset V$ and a single critical point.  Its
\emph{filled Julia set} is $K(g)=\{z:\ g^n(z)\in U\text{ for all }n\ge0\}$, and its
\emph{fundamental annulus} is $V\setminus\overline U$, with modulus
$\modulus(V\setminus\overline U)$.  A quadratic polynomial $f$ is
\emph{$p$-renormalizable} if $f^p\colon U\to V$ is quadratic-like with connected Julia
set, for disks $0\in U\Subset V$; the map $f^p$ is then a \emph{$p$-renormalization}.
Its small filled Julia sets are $K_p(i)=f^i(K_p)$ for $0\le i<p$, where $K_p$ is the
central one, and $\mathcal K_p=\bigcup_{i<p}K_p(i)$.

\begin{definition}[Primitive renormalization and bounded primitive type]
\label{def:apriori-primitive}
A $p$-renormalization is \emph{primitive} if the small filled Julia sets
$K_p(0),\dots,K_p(p-1)$ are pairwise disjoint.  A polynomial with primitive
renormalizations of periods $p_0<p_1<\cdots$ has \emph{$B$-bounded primitive type} if
\[
\frac{p_{k+1}}{p_k}\ \le\ B\qquad\text{for every }k.
\]
\emph{Source:} Kahn I, \texttt{intro.tex}, lines 41--57.
\end{definition}

\begin{definition}[A priori bounds]
\label{def:series-apriori}
An infinitely renormalizable quadratic polynomial $f$ has \emph{a priori bounds} if there
are $\epsilon>0$ and quadratic-like renormalizations $f^{p_k}\colon U_k\to V_k$ with
\[
\modulus(V_k\setminus U_k)\ \ge\ \epsilon\qquad\text{for all }k.
\]
\emph{Source:} Kahn I, \texttt{intro.tex}, lines 58--60.  Kahn III and Kahn--Lyubich II
use the same notion with the annulus written as $U_n\setminus U_n'$.
\end{definition}

\begin{definition}[Beau a priori bounds]
\label{def:apriori-beau}
A family $\mathcal C$ of renormalization combinatorics has \emph{beau a priori bounds} if
there is $\epsilon=\epsilon(\mathcal C)>0$ and, for every $\delta>0$, an integer
$N(\delta)$ such that every renormalizable quadratic-like map in $\mathcal C$ whose
initial modulus is at least $\delta$ has all renormalizations of level at least
$N(\delta)$ with modulus at least $\epsilon$.  The threshold $N(\delta)$ may depend on the
initial modulus, while $\epsilon$ depends only on the family.  \emph{Sources:}
Kahn III, \texttt{moduli.tex}, lines 623--633; Dudko--Lyubich (Feigenbaum),
\texttt{Feigen.tex}, lines 320--322; KKL, \texttt{waves-with-Misha.tex}, lines 117--123.
\end{definition}

\paragraph{The rigidity input.}
The a priori bounds sources of this chapter (Kahn I--III, Kahn--Kapiamba--Lyubich, and
Dudko--Lyubich at Feigenbaum points) pass from bounds to local connectivity through one cited
input, the rigidity theorem of Lyubich (Acta Math.~178 (1997)).  The sources state it in different
forms, and the book keeps those forms apart.  The satellite paper (Section~\ref{sec:dl-satellite})
uses a related input for Julia sets, and the limbs paper (Section~\ref{sec:dudko-schleicher})
does not use it.  Kahn I says that a priori bounds of
bounded primitive type give ``the Mandelbrot set is locally connected at $c$, or
equivalently $f_c$ is combinatorially rigid'' (\texttt{intro.tex}, lines 93--101), citing
Lyubich's [puzzle] entry.  Kahn II and III state the same implication for the decoration
and molecule conditions (Corollary of Kahn II, \texttt{intro.tex}, lines 123--126;
Corollary~6.4 of Kahn III, \texttt{intro.tex}, lines 36--41).  Dudko--Lyubich deduce MLC
at Feigenbaum points from beau bounds and the Rigidity Theorem of Lyubich.  KKL show that
beau bounds alone do not give MLC: rigidity additionally needs a bounded puzzle
Teichm\"uller diameter, and KKL state that Lyubich's argument needs only this weaker
hypothesis (KKL, Theorem~B).  Throughout this chapter, ``rigidity input'' means this
cited principle, and no digest below proves it.

\section{Kahn I: bounded primitive combinatorics}
\label{sec:kahn-I}

\paragraph{Identity.}
J.~Kahn, \emph{A priori bounds for some infinitely renormalizable quadratics: I.
Bounded primitive combinatorics}, arXiv:math/0609045 (version~3, 2006) \cite{KahnPrimitive}.  The title and
the single author are confirmed by the downloaded metadata and by the main TeX file.

\paragraph{Definitions used.}
\begin{itemize}
\item \emph{Geodesic lengths.} For a primitive renormalization with small Julia sets
$\mathcal K_p$, $\gamma_p$ is the closed hyperbolic geodesic surrounding $K_p$ in
$\C\setminus\mathcal K_p$.  For a quadratic-like $g\colon U\to V$, $\gamma(g)$ is the
core geodesic of $V\setminus K(g)$, and $|\eta|$ is hyperbolic length.
(\texttt{intro.tex}, lines 45--51.)
\item \emph{Pseudo-quadratic-like maps.} A pair $(i,f)\colon U'\to U$ in which $i$ is an
immersion, $f$ is a branched covering of degree $2$, $i^{-1}(K)=f^{-1}(K)=K'$ and $K'$ is
connected.  They carry canonical renormalizations, which is why the proof works with them;
they are introduced in this paper and cited in later sources as ``$\psi$-ql maps''
(\texttt{psiql2.tex}, lines 15--34).
\item \emph{Canonical weighted arc diagram.} $W_{\rm can}=\sum_\alpha W(\alpha)\,\alpha$,
where arcs cross thin rectangles and $W(\alpha)$ is their conformal width.  Horizontal arcs
join components of $\mathcal K_p$; vertical arcs join $\mathcal K_p$ to $\partial V$
(\texttt{intro.tex}, lines 229--245).
\end{itemize}

\paragraph{Results.}

\begin{theorem}[Length comparison]
\label{thm:kahn-I-length}
There is an absolute constant $\epsilon>0$ such that for every $p$ there is $M(p)$ with the
following property.  If $g\colon U\to V$ is a primitively $p$-renormalizable quadratic-like
map, then
\[
|\gamma_p|>M(p)\qquad\Longrightarrow\qquad |\gamma(g)|\ \ge\ \epsilon\,p\,|\gamma_p|.
\]
\emph{Status:} proved in the paper (source label \texttt{thm:length-comparison};
\texttt{intro.tex}, lines 159--167; \texttt{conclusions.tex}, lines 67--69).
\end{theorem}

\begin{theorem}[Bad now, worse earlier]
\label{thm:kahn-I-bnwe}
There is $l\in\N$ such that for every $B$ there is $M$ with the following property.  For
$f$ satisfying the standard assumptions of bounded primitive type $B$, and for all $n$,
\[
|\gamma_{(n+l)}|>M\qquad\Longrightarrow\qquad |\gamma_{(n)}|>2\,|\gamma_{(n+l)}|.
\]
\emph{Status:} proved in the paper (source label \texttt{thm:bnwe};
\texttt{conclusions.tex}, lines 137--145).
\end{theorem}

\begin{remark}
\label{rem:kahn-I-improving}
The bounded-period improvement used by Kahn II and Kahn III is a result of Kahn I.  In the
present version of Kahn I it appears as the ``improving of length'' theorem
(\texttt{conclusions.tex}, lines 74--82 and 172); Kahn II and Kahn III cite it as Theorem~9.1
of the first arXiv version.  Its statement is recorded in
Theorem~\ref{thm:kahn-II-bounded-period}.
\end{remark}

\begin{theorem}[Geodesic-to-modulus]
\label{thm:kahn-I-geodesic-modulus}
If $f\colon V'\to V$ is primitively $p$-renormalizable and $|\gamma_p|\le M$, then its
Douady--Hubbard $p$-renormalization has modulus at least $m(M)>0$.  \emph{Status:} proved
in the paper, using the canonical renormalization.  For polynomials the statement is McMullen's Theorem~10.3 in
\emph{Complex dynamics and renormalization} \cite{McMullenCDR}, for $M$ sufficiently small, and Theorem~4.10 in
\emph{Renormalization and 3-manifolds which fiber over the circle} \cite{McMullenTowers}, in general, for a slightly
different statement that readily implies this one (\texttt{conclusions.tex}, lines 246--258).
\end{theorem}

\begin{theorem}[Beau bounds for pseudo-quadratic-like maps]
\label{thm:kahn-I-psi-beau}
For every $B$ there is $M_0$, and for every $m$ there are $n_1$ and $M_1$, such that the following holds.  If $f$
satisfies the standard assumptions of bounded primitive type $B$ and $|\gamma(f)|\le m$, then $|\gamma_{(n)}|\le M_1$
for all $n\in\N$, and $|\gamma_{(n)}|\le M_0$ for all $n\ge n_1$ (\texttt{conclusions.tex}, lines 186--198).
\emph{Status:} proved (source label \texttt{thm:psi-beau};
\texttt{conclusions.tex}, lines 186--198).
\end{theorem}

\begin{theorem}[Beau bounds for semi-quadratic-like maps]
\label{thm:kahn-I-beau-semi}
For all $B$ there is $M_0$ such that for all $m$ there are $n_1$ and $M_1$ with the following property.  Let
$f\colon U'\to U$ be a semi-quadratic-like map satisfying the standard assumptions, with $|\modulus(V\setminus K)|\ge m$.
Then (1) for all $n\in\N$ there is a quadratic-like $p_n$-renormalization $f^{p_n}\colon V'\to V$ with
$\modulus(V\setminus V')\ge M_1$, and (2) for $n\ge n_1$ there is one with $\modulus(V\setminus V')\ge M_0$.
\emph{Status:} proved in the paper, as an immediate consequence of the pseudo-quadratic-like beau bounds and of the
geodesic-to-modulus theorem (source label \texttt{thm:a-priori-bounds}; \texttt{conclusions.tex}, lines 279--296).
\end{theorem}

\begin{theorem}[Main Theorem, Kahn I]
\label{thm:kahn-I-main}
If $f$ is an infinitely renormalizable quadratic polynomial of bounded primitive type, then
$f$ has a priori bounds.  \emph{Status:} proved in the paper (\texttt{intro.tex}, lines
69--74).  The explicit form is the following.  For every $B$
there is $m>0$ such that every infinitely $B$-bounded primitively renormalizable quadratic
polynomial has, at every primitive period $p_n$, a quadratic-like renormalization $V\supset V'$
with $\modulus(V\setminus V')\ge m$ (the explicit form is stated without a label in the source;
\texttt{conclusions.tex}, lines 298--304).
\end{theorem}

\paragraph{Proof strategy (from the source).}
Modulus bounds are reduced to bounded lengths of the geodesics $\gamma_{(n)}$
(\texttt{intro.tex}, lines 128--139).  The argument localizes between periods $p_{n-l}$ and
$p_n$; bounded period ratios give at most $B^l$ relevant small Julia sets.  Pseudo-quadratic-like
maps and canonical renormalization make the relevant lengths compare monotonically.  Canonical
weighted arc diagrams and extremal width estimates control the vertical width, and horizontal
weight controls vertical weight through the primitive combinatorics and an entropy theorem
(\texttt{intro.tex}, lines 141--245; \texttt{wad.tex}, lines 83--162).

\paragraph{External inputs.}
Lyubich, \emph{Dynamics of quadratic polynomials I--II} (rigidity input, source key
\texttt{puzzle}); McMullen, \emph{Complex dynamics and renormalization} (source key
\texttt{McM}); McMullen, \emph{Renormalization and 3-manifolds which fiber over the circle}
(source key \texttt{towers}); Sullivan (source key \texttt{S}, real bounded-type a priori
bounds and the term ``beau''); Yoccoz (source key \texttt{H}, local connectivity of
non-infinitely renormalizable Julia sets); Hu--Jiang and Jiang (source keys \texttt{HJ},
\texttt{J}, local connectivity of $J(f_c)$).

\paragraph{MLC scope.}
The paper lists as a consequence of its a priori bounds that $\M$ is locally connected at the parameters it covers, namely the
infinitely renormalizable $c$ of bounded primitive type, equivalently combinatorial rigidity
of $f_c$, the implication being attributed to Lyubich's rigidity theorem (\texttt{intro.tex}, lines 93--101).  The introduction also says that the class of ``high bounded type'' had been treated before, citing the same reference.  It does not claim MLC globally.  It also
lists further consequences (local connectivity of $J(f_c)$ and universality statements),
which we do not digest further (\texttt{intro.tex}, lines 102--119).

\section{Kahn II: decorations}
\label{sec:kahn-II}

\paragraph{Identity.}
J.~Kahn and M.~Lyubich, \emph{A priori bounds for some infinitely renormalizable
quadratics: II. Decorations}, arXiv:math/0609046 (2006) \cite{KahnLyubichDecorations}.  The main TeX file is
\texttt{refs/arxiv-math-0609046/main.tex}, lines 39--41.

\paragraph{Definitions.}
\begin{itemize}
\item \emph{Decoration.} A decoration $\mathcal L$ of $\M$ is the part of $\M$ cut off by two
external rays landing at a tip of a satellite copy of $\M$ attached to the main cardioid
(\texttt{main.tex}, lines 49--50; \texttt{intro.tex}, lines 88--91).
\item \emph{Escape route.} For a map that is not satellite renormalizable, let $n$ be least
with $f^{qn}(0)\in\operatorname{int}Z^1_{\boldsymbol\kappa}$.  The escape route
$\bar\kappa=(\kappa_1,\dots,\kappa_n=\boldsymbol\kappa)$ satisfies
\[
f^{qm}(0)\in Z^{\,q(n-m)+1}_{\boldsymbol\kappa_m},\qquad m=1,\dots,n
\]
(\texttt{kingdom.tex}, lines 337--353).
\item \emph{Decoration condition.} An infinitely renormalizable $f$ satisfies the decoration
condition if all little copies $M^{(n)}$ belong to finitely many decorations
$\mathcal L_k$; equivalently, the parameters $q_n$ and $n_n$ of the escape routes are bounded
(\texttt{kingdom.tex}, lines 501--504).
\item \emph{Decoration class $(\bar q,\bar n)$.} A $\psi$-ql map belongs to it when every
renormalization has decoration parameters $(\bar q_n,\bar n_n)\le(\bar q,\bar n)$
(\texttt{moduli.tex}, lines 1335--1337).
\end{itemize}

\paragraph{Results.}

\begin{theorem}[Main Theorem, Kahn II]
\label{thm:kahn-II-main}
An infinitely renormalizable quadratic map satisfying the decoration condition has a priori
bounds.  \emph{Status:} proved (\texttt{intro.tex}, lines 112--118).  The proof uses the
bounded-period improvement (Theorem~\ref{thm:kahn-II-bounded-period}, which is Kahn I's result)
for the regime of bounded periods.
\end{theorem}

\begin{theorem}[High-period theorem]
\label{thm:kahn-II-high}
Given finitely many decorations $\mathcal L_k$ there is $\underline p$ such that every
infinitely renormalizable quadratic map satisfying the decoration condition with these
decorations and renormalization periods $p\ge\underline p$ has a priori bounds.
\emph{Status:} proved in the paper (\texttt{intro.tex}, lines 132--136).
\end{theorem}

\begin{theorem}[Improving the moduli, bounded decoration parameters]
\label{thm:kahn-II-moduli}
For every $(\bar q,\bar n)$ and $\rho>0$ there are $\underline p\in\N$ and $\epsilon>0$ such
that the following holds.  Let $f\colon(U,K)\to(U,K)$ be a renormalizable $\psi$-ql map with
decoration parameters at most $(\bar q,\bar n)$ and first renormalization $f'=f^p\colon(U',K')\to(U',K')$.
If $p\ge\underline p$ and $\modulus(U'\setminus K')<\epsilon$, then
\[
\modulus(U\setminus K)\ <\ \rho\,\modulus(U'\setminus K').
\]
\emph{Status:} proved (\texttt{moduli.tex}, lines 1271--1281).
\end{theorem}

\begin{theorem}[Improving the moduli, bounded period]
\label{thm:kahn-II-bounded-period}
For $\rho\in(0,1)$ there is $\underline p(\rho)$ such that for every $\bar p\ge\underline p$
there is $\epsilon(\bar p)>0$ with the following property.  If $f$ is primitively
renormalizable, $\underline p\le p\le\bar p$ and $\modulus(U'\setminus K')<\epsilon(\bar p)$,
then $\modulus(U\setminus K)<\rho\,\modulus(U'\setminus K')$.
\emph{Status:} \emph{cited}, not proved in Kahn II.  The source identifies it as Theorem~9.1
of Kahn I, which is cited as it appeared in the first arXiv version of Kahn I
(\texttt{moduli.tex}, lines 1309--1327; Kahn I, \texttt{conclusions.tex}, lines 74--82).
Kahn I's present version contains this result as the ``improving of length'' theorem
(\texttt{conclusions.tex}, line 172).  We record the dependency and not the numbering.
\end{theorem}

\begin{theorem}[Corollary: geometric decay]
\label{thm:kahn-II-decay}
For each $(\bar q,\bar n)$ there are $\epsilon>0$ and $l\in\N$ such that for every
infinitely renormalizable $\psi$-ql map in the decoration class $(\bar q,\bar n)$, if
$\modulus(U_n\setminus K_n)<\epsilon$ for some $n\ge l$, then
\[
\modulus(U_{n-l}\setminus K_{n-l})\ <\ \tfrac12\,\modulus(U_n\setminus K_n).
\]
\emph{Status:} proved (\texttt{moduli.tex}, lines 1340--1346).
\end{theorem}

\begin{theorem}[Beau bounds in the decoration class]
\label{thm:kahn-II-beau}
For every $(\bar q,\bar n)$ the family of renormalization combinatorics in the decoration class
$(\bar q,\bar n)$ has beau a priori bounds (Definition~\ref{def:apriori-beau}).
\emph{Status:} proved (\texttt{moduli.tex}, lines 1423--1426).
\end{theorem}

\begin{theorem}[JLC and MLC implication]
\label{thm:kahn-II-mlc}
If $f_c(z)=z^2+c$ is infinitely renormalizable and satisfies the decoration condition, then
$J(f_c)$ is locally connected and $\M$ is locally connected at $c$.
\emph{Status:} the implication is cited from Lyubich (rigidity input,
Section~\ref{sec:apriori-framework}); the paper does not prove it
(\texttt{intro.tex}, lines 36--37 and 123--126).
\end{theorem}

\paragraph{Proof strategy (from the source).}
The combinatorics are encoded by the Yoccoz puzzle and a modified principal nest
(\texttt{intro.tex}, lines 155--166).  Ordinary quadratic-like maps are replaced by canonical
$\psi$-ql maps with pseudo-puzzles (\texttt{intro.tex}, lines 168--179).  The Quasi-Additivity Law
and the Covering Lemma supply the analytic estimates (\texttt{intro.tex}, lines 181--183).
Modulus improvement for bounded decoration parameters is proved first
(\texttt{moduli.tex}, lines 1271--1307), then combined with the bounded-period theorem to cover all
period regimes (\texttt{moduli.tex}, lines 1309--1407).  Iteration gives beau bounds, which are
transferred back to ordinary quadratic-like maps (\texttt{moduli.tex}, lines 1409--1440).

\paragraph{External inputs.}
Lyubich, \emph{Dynamics of quadratic polynomials I--II} (the implication above); Kahn I (the
bounded-period theorem); Kahn--Lyubich, \emph{The quasi-additivity law in conformal geometry}
(analytic tools); Kahn--Lyubich, \emph{Local connectivity of Julia sets for unicritical
polynomials} (the modified principal nest); Douady--Hubbard, \emph{On the dynamics of
polynomial-like maps} (background).  The transfer from $\psi$-ql to ordinary quadratic-like
maps uses a ``ql extension'' lemma; the source does not say where that lemma comes from
(\texttt{moduli.tex}, lines 1439--1440).

\paragraph{MLC scope.}
The claimed scope is quadratic, infinitely renormalizable parameters satisfying the decoration
condition: local connectivity of $J(f_c)$ and of $\M$ at $c$ (\texttt{intro.tex}, lines 123--126).
The source also states that its a priori bounds proof extends unchanged to higher-degree
unicritical maps, but that the MLC implication uses quadratic geometry; for higher degree, the
source cites Cheraghi together with Avila--Kahn--Lyubich--Shen (\texttt{intro.tex}, lines
199--203).  It does not claim global local connectivity of $\M$.

\section{Kahn III: molecules}
\label{sec:kahn-III}

\paragraph{Identity.}
J.~Kahn and M.~Lyubich, \emph{A priori bounds for some infinitely renormalizable quadratics:
III. Molecules}, arXiv:0712.2444 (2007) \cite{KahnLyubichMolecules}.  The main TeX file is \texttt{refs/arxiv-0712.2444/main.tex},
lines 20--22.

\paragraph{Definitions.}
\begin{itemize}
\item \emph{Molecule.} $\mathcal M$ is the closure of the union of all hyperbolic components
obtained by iteratively attaching secondary hyperbolic components to the main cardioid; the
source also calls it the cactus (\texttt{intro.tex}, lines 3--11).  Combinatorially (\texttt{kingdom.tex}, lines 169--176), $\mathcal M$ consists of the maps that
are either finitely renormalizable, with every renormalization satellite and the last one having a
non-repelling cycle, or infinitely renormalizable with every renormalization satellite.
\item \emph{$\eta$-molecule condition.} For primitively renormalizable maps $f_i$, the associated
copies $M$ stay $\eta$-away from $\mathcal M$.  For an infinitely primitively renormalizable $f$, the
relative copies $\widetilde M_i$ associated to $R^if$ stay $\eta$-away
(\texttt{kingdom.tex}, lines 195--204).
\item \emph{Parameters $(r,q,n)$.} $r$ is the period of a dividing cycle with an $r$-ray portrait;
$f^r$ has a repelling fixed point $\alpha$ with $q$ rays landing there; the orbit $f^{rj}(0)$ does
not escape the central domain for $j=1,\dots,qn-1$; $n$ is the first time $f^{rqn}(0)$ escapes it
(\texttt{kingdom.tex}, lines 208--240).
\item \emph{$(\bar r,\bar q,\bar n)$-molecule condition.} There is a choice of $(r,q,n)$ with
$r\le\bar r$, $q\le\bar q$, $n\le\bar n$.  The $(r,q,n)$ condition and the $\eta$-condition are
equivalent, with quantitative dependence in both directions (\texttt{kingdom.tex}, lines 242--251).
\end{itemize}

\paragraph{Results.}

\begin{theorem}[Main Theorem, Kahn III]
\label{thm:kahn-III-main}
Infinitely renormalizable quadratic maps satisfying the molecule condition have a priori bounds.
\emph{Status:} proved here, using Kahn I's bounded-period theorem (\texttt{intro.tex}, lines 28--34).
\end{theorem}

\begin{theorem}[High periods]
\label{thm:kahn-III-high}
Given $\eta>0$ and $\rho\in(0,1)$ there are $\bar\mu>0$ and $\underline p\in\N$ such that, if a
period-$p$ renormalizable quadratic-like $f$ satisfies the $\eta$-molecule condition,
$\operatorname{psimod}(Rf)<\bar\mu$ and $p\ge\underline p$, then
\[
\operatorname{psimod}(f)\ <\ \rho\,\operatorname{psimod}(Rf).
\]
Here $\operatorname{psimod}$ is the modulus of the $\psi$-ql map as defined in the source.
\emph{Status:} proved here (\texttt{intro.tex}, lines 51--60).
\end{theorem}

\begin{theorem}[Improving moduli for a fixed molecule class]
\label{thm:kahn-III-moduli}
For fixed $(\bar r,\bar q,\bar n)$ and $\rho>0$ there are $\underline p$ and $\epsilon>0$ such that
a renormalizable $\psi$-ql $f$ of period $p\ge\underline p$ satisfying the $(\bar r,\bar q,\bar n)$
condition has $\operatorname{psimod}(f)<\rho\,\modulus(g)$ whenever its first renormalization $g$
has $\operatorname{psimod}(g)<\epsilon$.  \emph{Status:} proved here (\texttt{moduli.tex}, lines 579--591).
\end{theorem}

\begin{theorem}[Bounded periods]
\label{thm:kahn-III-bounded}
The bounded-period improvement of Theorem~\ref{thm:kahn-II-bounded-period} is the input used here.
\emph{Status:} \emph{cited}, specifically Theorem~9.1 of Kahn I, as cited in Kahn III
(\texttt{moduli.tex}, lines 593--611).
\end{theorem}

\begin{theorem}[Corollary 6.3]
\label{thm:kahn-III-decay}
For fixed $(\bar r,\bar q,\bar n)$ there are $\epsilon>0$ and $l$ such that, for infinitely
renormalizable $f$ satisfying the condition, $\operatorname{psimod}(g_n)<\epsilon$ and $n\ge l$ imply
\[
\operatorname{psimod}(g_{n-l})\ <\ \tfrac12\,\operatorname{psimod}(g_n).
\]
\emph{Status:} proved here from the two preceding theorems (\texttt{moduli.tex}, lines 613--619).
\end{theorem}

\begin{theorem}[Beau bounds for molecule classes]
\label{thm:kahn-III-beau}
Every family of renormalization combinatorics satisfying a fixed $(\bar r,\bar q,\bar n)$-molecule
condition has beau a priori bounds.  \emph{Status:} proved here (\texttt{moduli.tex}, lines 635--640).
\end{theorem}

\begin{theorem}[Corollary 6.4: JLC and MLC]
\label{thm:kahn-III-mlc}
If $f_c(z)=z^2+c$ is infinitely renormalizable and satisfies the molecule condition, then $J(f_c)$ is
locally connected and $\M$ is locally connected at $c$.  \emph{Status:} conditional on the rigidity
input of Lyubich, \emph{Dynamics of quadratic polynomials I--II} (\texttt{intro.tex}, lines 36--41).
\end{theorem}

\paragraph{Proof strategy (from the source).}
Primitive renormalizations are encoded by nested relative $\M$-copies and bounded $(r,q,n)$ data
(\texttt{kingdom.tex}, lines 178--211).  The Yoccoz puzzle and principal nest, with $\psi$-moduli,
supply the geometry (\texttt{intro.tex}, lines 78--100).  The high-period improvement compares the
bottom and top principal-nest geometry (\texttt{moduli.tex}, lines 1--20).  It is combined with the
bounded-period theorem (\texttt{moduli.tex}, lines 593--622), and iterating the improvement gives the
uniform and beau bounds (\texttt{moduli.tex}, lines 613--640).

\paragraph{External inputs.}
Kahn I (bounded-period theorem); Kahn--Lyubich, \emph{Quasi-additivity law in conformal geometry}
(proof machinery); Kahn--Lyubich II (analogous results and Lemma~5.3); Lyubich, \emph{Dynamics of
quadratic polynomials I--II} (the implication for $J(f_c)$ and $\M$); Douady--Hubbard, \emph{On the
dynamics of polynomial-like maps} (background).

\paragraph{MLC scope.}
The paper claims local connectivity of $J(f_c)$ and of $\M$ at the corresponding parameter, for
infinitely renormalizable quadratics satisfying the molecule condition (\texttt{intro.tex}, lines
30--41).  It does not claim global local connectivity of $\M$, nor local connectivity for all
infinitely renormalizable quadratics (\texttt{intro.tex}, lines 12--15 and 43--48).

\section{Kahn IV: elephant eyes}
\label{sec:kahn-IV}

\paragraph{Identity.}
J.~Kahn and M.~Lyubich, \emph{A priori bounds for some infinitely renormalizable quadratics:
IV. Elephant Eyes}, arXiv:2601.21905 (2026) \cite{KahnLyubichElephant}.  The main TeX file is
\texttt{refs/arxiv-2601.21905/main.tex}.  The source says that the result was obtained around
2010 and reproduced from a January 2020 version (\texttt{intro.tex}, lines 3--13).

\paragraph{Scope.}
Parts I--III treat bounded combinatorics.  Part IV treats unbounded combinatorics: arbitrary
$1/q$-limbs with renormalization periods $p=q+O(1)$ (\texttt{intro.tex}, lines 18--29).

\paragraph{Definitions.}
\begin{itemize}
\item \emph{Setup.} Start with an even quadratic-like map $f\colon U'\to U$ with connected Julia set
and both fixed points repelling.  The fixed point $\beta$ is the landing point of the external
$0$-ray; $\alpha$ is the landing point of $q>1$ rays, cyclically permuted with rotation number
$p/q$.  These rays define sectors $S_i$ (\texttt{elcomb.tex}, lines 5--20).
\item \emph{Elephant-eye combinatorics} $(\mathrm E1)$--$(\mathrm E3)$.  $(\mathrm E1)$: $\alpha$ has
rotation number $1/q$; $\Upsilon'$ maps two-to-one onto $S_1$; each $S_i'$ and $Z_i'$ maps univalently
onto $S_{i+1}$ for $i=1,\dots,q-2$; the terminal sectors map univalently onto
$\Upsilon\cup\bigcup_iZ_i\cup(-R_i)$ (\texttt{elcomb.tex}, lines 36--47).  $(\mathrm E2)$: $K_0$ is the
only little Julia set in $\Upsilon$ (lines 49--50).  $(\mathrm E3)$: $p=q+b$ with $b\le\bar b$ bounded,
$q$ unbounded, and, for simplicity, $b<q$ (lines 52--59).
\item \emph{Hubbard tree of little Julia sets.}  The tree spanned by the little Julia sets, collapsed to
points; its vertices are the collapsed $K_i$, the fixed point $\alpha$, and branch points
(\texttt{elcomb.tex}, lines 61--73).
\item \emph{A priori bounds.}  The source gives no single formal definition.  Its geometric content is a
fundamental annulus of definite modulus which avoids the non-central top-level little Julia sets
(``unbranched''; \texttt{lc.tex}, lines 2--8).
\end{itemize}

\paragraph{Results.}

\begin{theorem}[Main estimate]
\label{thm:kahn-IV-main-estimate}
For every ideal marking $\mathcal I=(I_1,\dots,I_p)$ of the hyperbolic plane $\D$,
\[
\mathcal W(\D,\mathcal I)-O(p)\ \le\ 2\sum_{\{I_m,I_n\}}\mathcal W\{I_m,I_n\}\ \le\ \mathcal W(\D,\mathcal I),
\]
where the sum runs over unordered pairs.  \emph{Status:} proved in the paper
(\texttt{canon.tex}, lines 172--183).  This is the disk case of the thin--thick decomposition that
the abstract calls the new tool for bordered Riemann surfaces (\texttt{main.tex}, lines 34--40).
\end{theorem}

\begin{theorem}[Bounded degree]
\label{thm:kahn-IV-bounded-degree}
Every Hubbard-tree edge maps under $f^p$ at most two-to-one onto at most $b+1$ consecutive edges.
\emph{Status:} proved in the paper (\texttt{elcomb.tex}, lines 168--190).
\end{theorem}

\begin{theorem}[Amplification]
\label{thm:kahn-IV-amplification}
For every $\lambda>1$ there are thresholds $\underline W_{\rm loc}(N)$ and $\underline p(\lambda,N,L)$ such
that, for sufficiently large $N$ and $L$, the vertical weight satisfies $W^{\rm ver}\ge\lambda\,W_{\rm loc}$
whenever $W_{\rm loc}\ge\underline W_{\rm loc}(N)$ and $p\ge\underline p(\lambda,N,L)$.
\emph{Status:} proved in the paper (\texttt{confined.tex}, lines 180--196).
\end{theorem}

\begin{remark}
\label{rem:kahn-IV-unclear}
The macro file declares two environments, ``Main Theorem (physical version)'' and ``Main Theorem
(dynamical version)'', but never states them (\texttt{macros.tex}, lines 17--18).  So the final
a priori bounds statement and its constants are not recoverable as a numbered theorem from this
TeX source.  We record only the estimates above, and the abstract's claim that the paper proves a
priori bounds for elephant-eye combinatorics.
\end{remark}

\paragraph{Local connectivity and MLC claims.}
The source states local connectivity of the Julia sets under consideration by citing Hu--Jiang and Jiang,
and MLC at the corresponding parameters by Lyubich's machinery together with little Mandelbrot copies on
Ecalle--Voronin cylinders (\texttt{lc.tex}, lines 2--17).  The external inputs are Lyubich 1997, \S\S 9.2 and 11
(\texttt{lc.tex}, lines 9--11 and 46--50).  The source does not cite an unpublished ``Elephant puzzles''
manuscript; the unpublished inputs it does cite are Dudko--Lyubich--Peterson, \emph{Taming the elephants}
(in preparation) and Cheritat's work on small copies in elephants (\texttt{bib.tex}, lines 14--24).

\paragraph{External inputs.}
Lyubich, \emph{Dynamics of quadratic polynomials I--II} (\texttt{bib.tex}, lines 66--67); Kahn I, II and III
(\texttt{bib.tex}, lines 38--59); Hu--Jiang (unpublished, 1993) and Jiang (\texttt{bib.tex}, lines 26--32);
Douady and Lavaurs, geometric limits of parabolic bifurcations (\texttt{bib.tex}, lines 18--20 and 34--36).

\paragraph{MLC scope.}
The claim is local connectivity and MLC for elephant-eye combinatorics, with $q$ unbounded and
$p=q+O(1)$.  It is not a claim for all infinitely renormalizable parameters, and the paper does not claim
the full MLC conjecture (\texttt{intro.tex}, lines 18--29; \texttt{lc.tex}, lines 9--17).

\paragraph{Proof strategy (from the source).}
The combinatorics are encoded by Hubbard-tree arcs and a bounded-degree lemma for $f^p$
(\texttt{elcomb.tex}, lines 129--190).  Canonical laminations and thin--thick decompositions are
developed for marked hyperbolic planes and bordered Riemann surfaces (\texttt{canon.tex}, lines 150--183
and 275--520).  Geometric widths are converted into dynamical weights (\texttt{weights.tex}, lines 1--20 and
228--250), pull-off arguments handle bounded and unbounded combinatorics (\texttt{pulloff.tex}, lines 1--12;
\texttt{confined.tex}, lines 1--14), and vertical weight is amplified (\texttt{confined.tex}, lines 180--196 and
582--712).  Unbranched a priori bounds then feed the rigidity and local connectivity machinery
(\texttt{lc.tex}, lines 2--18).

\section{Dudko--Lyubich: satellite parameters}
\label{sec:dl-satellite}

\paragraph{Identity.}
D.~Dudko and M.~Lyubich, \emph{Local connectivity of the Mandelbrot set at some satellite parameters of
bounded type}, arXiv:1808.10425; Geometric and Functional Analysis 33 (2023), no.~4, 912--1047
\cite{DudkoLyubichSatellite}.  The main file is \texttt{refs/arxiv-1808.10425/MLCatSatParam.tex}, whose title
and authors are at lines 713--718.  The Feigenbaum paper cites this work under the title
``parameter values of bounded type''; we use the arXiv title.

\paragraph{Definitions.}
\begin{itemize}
\item \emph{Infinitely renormalizable parameter.}  $c\in\bigcap_{i\ge0}\operatorname{Dom}\mathcal R_{\rm QL}^i$
(\texttt{MLCatSatParam.tex}, lines 2238--2244).
\item \emph{Bounded (satellite) type.}  The $\mathcal R_{\rm QL}$-orbit of $c$ lies in finitely many maximal
small Mandelbrot copies; the type is \emph{bounded satellite type} if all these copies are satellite copies
(lines 2243--2247).
\item \emph{Primary satellite copies.}  For $r\in\mathbb Q\cap(0,1)$, $\mathcal M_r$ is the unique copy attached at
the parabolic parameter $c(r)$ on the main cardioid with rotation number $r$ (lines 2250--2254).  Iterated copies
are $\mathcal M_{r,s}:=\mathcal R_r^{-1}(\mathcal M_s)$ and $\mathcal M_{r,s,t}:=\mathcal R_{r,s}^{-1}(\mathcal M_t)$
(lines 2254--2259).
\item \emph{Construction.}  Fix an $\mathcal R_c$-periodic $c(\theta)\in\partial\mathcal H$ of period $m$ and a copy
$\mathcal M_0$ near $c(\theta)$; $\mathcal M_{n-1}$ is the unique inverse image under $\mathcal R_c^m$ still near
$c(\theta)$, so $\mathcal M_n$ shrinks to $c(\theta)$ (lines 876--893).  The associated parameter is $c_n$
(lines 898--905).
\item \emph{Unbranched a priori bounds.}  For infinitely many renormalizations $f_n\colon U_n\to V_n$,
$V_n\cap\operatorname{Post}(f)\subset K(f_n)$ and $\operatorname{mod}(V_n\setminus U_n)\ge\varepsilon>0$
(lines 2261--2269).
\item \emph{Valuable flower.}  The combinatorial connected hull of the cycle of secondary small filled Julia sets;
it is forward invariant and contains the postcritical set (lines 998--1002 and 7084--7095).
\end{itemize}

\paragraph{Results.}

\begin{theorem}[Main theorem of the satellite paper]
\label{thm:dl-sat-main}
There is a small copy $\mathcal M_0$ near $c(\theta)$ such that the map
\[
\mathcal R_c^m\colon\ \{c(\theta)\}\cup\bigcup_{n<0}\mathcal M_n\ \longrightarrow\
\{c(\theta)\}\cup\bigcup_{n\le0}\mathcal M_n
\]
is $(1+\varepsilon)$-conformal at $c(\theta)$.  For every $n\le0$, $\bigcap_{i\ge0}\operatorname{Dom}(\mathcal R_n^i)=\{c_n\}$,
the Julia set $J_{c_n}$ is locally connected, and $J_{c_n}$ has positive area.
\emph{Status:} proved in the paper (\texttt{MLCatSatParam.tex}, lines 894--907).  The positive-area claim
uses the Avila--Lyubich estimates, which are cited (lines 1017--1023).
\end{theorem}

\begin{theorem}[Valuable flowers]
\label{thm:dl-sat-flowers}
The quadratic polynomial $p_n(z)=z^2+c_n$ has a forward-invariant valuable flower $X_n$ containing its
postcritical set and lying near the closed Siegel disk of $z^2+c(\theta)$ (\texttt{MLCatSatParam.tex}, lines 908--911).
\emph{Status:} proved in the paper.
\end{theorem}

\begin{theorem}[Ternary small copy]
\label{thm:dl-sat-ternary}
Near a selected ternary parabolic prepacman $\mathbf F_{r,s,f}$ there is a ternary satellite copy
$\mathbf M_0=\mathbf M_{r,s,k}$ (\texttt{MLCatSatParam.tex}, lines 6798--6823).  \emph{Status:} proved in the paper.
\end{theorem}

\begin{theorem}[Projection of valuable flowers]
\label{thm:dl-sat-projection}
For $n\ll0$ and every $\mathbf G_n\in\mathbf M_n$, the enlarged and ordinary valuable flowers project to nice
flowers in a prescribed neighbourhood of the Siegel disk (\texttt{MLCatSatParam.tex}, lines 7084--7095).
\emph{Status:} proved in the paper.
\end{theorem}

\paragraph{What is not proved.}
The paper does not prove a general scaling law for primary satellite copies; the authors say that this
remains suggested by pictures (lines 992--995 and 6818--6828).  It also says that the first partial copy
$\mathbf M(\mathcal H_{r,s})$ is not known to be bounded or complete (lines 6818--6828).

\paragraph{Proof strategy (from the source).}
A hyperbolic pacman-renormalization unstable manifold is built near a periodic Siegel map (lines 1013--1015).
A ternary satellite copy is built there from external rays, quadratic-like thickening, straightening and
Yoccoz's inequality (lines 6798--6823).  Valuable flowers are built in the transcendental plane and projected
to pacman planes (lines 998--1002), and flowers and hybrid classes are transferred through stable laminations
to quadratic polynomials (lines 1008--1015).  Hyperbolicity and asymptotically conformal holonomy give scaling
(lines 1013--1017).  Rigidity comes from expansion, and local connectivity of $J$ from unbranched a priori
bounds; positive area uses the Avila--Lyubich estimates (lines 1017--1023).

\paragraph{External inputs.}
Dudko--Lyubich--Selinger, \emph{Pacman renormalization and self-similarity of the Mandelbrot set near Siegel
parameters} (\texttt{bib.tex}, lines 57--59); Douady--Hubbard straightening (\texttt{bib.tex}, lines 49--52);
Yoccoz's inequality; Cheraghi--Shishikura, \emph{Satellite renormalization of quadratic polynomials},
arXiv:1509.07843 (context only); Avila--Lyubich, \emph{Lebesgue measure of Feigenbaum Julia sets} (2022), for
positive area; Jiang, McMullen and Lyubich, for local connectivity from unbranched a priori bounds; Lomonaco--Petersen,
for the obstruction to quasiconformal equivalence of satellite copies (\texttt{MLCatSatParam.tex}, lines 6824--6828).

\paragraph{MLC scope.}
The positive claim is local: $\M$ is locally connected at the constructed infinitely renormalizable parameters
$c_n$ of bounded satellite type, and $J_{c_n}$ is locally connected of positive area (lines 724--730 and 894--907).
The paper recalls that global MLC would require every connected component of the non-escaping set of
$\mathcal R_{\rm QL}$ to be a singleton (lines 2240--2244), so it is not a proof of global MLC.

\paragraph{Unclear.}
The main TeX source has no formal definition of an ``itinerary''; the paper uses landing-time indexing $P(i)$ of
puzzle pieces instead (lines 7095--7103).

\section{MLC at Feigenbaum points (Dudko--Lyubich)}
\label{sec:dl-feigenbaum}

\paragraph{Identity.}
D.~Dudko and M.~Lyubich, \emph{MLC at Feigenbaum points}, arXiv:2309.02107 (2023); Publications
Math\'ematiques de l'IH\'ES 143 (2026), 97--141 \cite{DudkoLyubichFeigenbaum}.  The main file is
\texttt{refs/arxiv-2309.02107/Feigen.tex}, whose title and authors are at lines 140--142.

\paragraph{Definitions.}
\begin{itemize}
\item \emph{Quadratic-like map and fundamental annulus.}  A holomorphic double branched covering
$f\colon X\to Y$ between Jordan disks $X\Subset Y\subset\C$, with one critical point (usually $0$).  The fundamental
annulus is $A=Y\setminus\overline X$, and $\modulus f:=\modulus A$ (\texttt{Feigen.tex}, lines 1679--1686).
\item \emph{Filled Julia set and germ.}  $K(f)=\{z:\ f^n(z)\in X\ \text{for all}\ n\ge0\}$; two maps with
connected Julia sets have the same germ if they share the filled Julia set and agree near it
(lines 1686--1700).
\item \emph{Hybrid conjugacy.}  A quasiconformal conjugacy $h\colon(Y,X)\to(\widetilde Y,\widetilde X)$ with
$\bar\partial h=0$ almost everywhere on the filled Julia set (lines 1706--1711).
\item \emph{Feigenbaum parameter.}  The paper calls the infinitely renormalizable quadratic polynomials of bounded type
Feigenbaum parameters.  The abstract makes this explicit, and the introduction extends the name to ``all infinitely
renormalizable parameters of bounded type'' (\texttt{Feigen.tex}, lines 150--165 and 210--216).  A Feigenbaum map is
one with all renormalization periods bounded by some $\bar p$ (lines 308--315).
\item \emph{Beau bounds.}  The constant $\mu>0$ depends only on $\bar p$, once $n$ is large; the threshold may
depend on the initial modulus (lines 315--322).
\item \emph{Renormalization conjecture.}  For every $\bar p$, $\mathcal R$ has a renormalization horseshoe
$\mathcal A_{\bar p}$ of type-$\bar p$ Feigenbaum maps on which it acts hyperbolically with one-dimensional
unstable foliation (lines 337--342).
\item \emph{Self-similarity over $\mathcal I_{\bar p}$.}  The straightening map
$\chi\colon\mathcal M^{n+1}(c)\to\mathcal M^n(\chi(c))$ is $C^{1+\delta}$-conformal at every $c\in\mathcal I_{\bar p}$,
uniformly (lines 346--355).
\end{itemize}

\paragraph{Results.}

\begin{theorem}[Main A]
\label{thm:dl-fei-A}
Every Feigenbaum quadratic-like map has a priori beau bounds.  \emph{Status:} proved in the paper
(\texttt{Feigen.tex}, lines 324--328).
\end{theorem}

\begin{theorem}[Main B]
\label{thm:dl-fei-B}
The Mandelbrot set is locally connected at every Feigenbaum parameter (\texttt{Feigen.tex}, lines 330--335).
\emph{Status:} deduced from Main A and Lyubich's Rigidity Theorem (\emph{Dynamics of quadratic polynomials I--II}, cited
in the paper as $[L_1]$; lines 3370--3374).  It is not an independent proof from scratch.
\end{theorem}

\begin{theorem}[Theorem C: renormalization conjecture]
\label{thm:dl-fei-C}
For every combinatorial bound $\bar p$, the Renormalization Conjecture is valid in the space of quadratic-like
maps (\texttt{Feigen.tex}, lines 337--344).  \emph{Status:} obtained in the paper from the cited universality theorem
of Lyubich (\emph{Feigenbaum--Coullet--Tresser universality and Milnor's hairiness conjecture}, Ann.\ of Math.\ 149
(1999), cited as $[L_2]$) together with Sullivan, McMullen, and Avila--Lyubich (lines 337--344 and 3376--3410).  We
record it as conditional on these inputs.
\end{theorem}

\begin{theorem}[Theorem D: self-similarity]
\label{thm:dl-fei-D}
For every $\bar p$, the Mandelbrot set is self-similar over $\mathcal I_{\bar p}$ (\texttt{Feigen.tex}, lines 346--360).
\emph{Status:} deduced from Theorem C; it is not a separate direct theorem.
\end{theorem}

\begin{proposition}[Teichm\"uller contraction]
\label{prop:dl-fei-teich}
For every $\bar p$ there is $\Delta_{\bar p}>1$ such that an infinitely renormalizable quadratic polynomial of bounded
type $\bar p$ has prerenormalizations $g_n=f_c^{p_n}\colon X_n\to Y_n$ with
\begin{itemize}
\item[(I)] $\operatorname{Width}(Y_n\setminus X_n)=O(\Delta_{\bar p}^{\,n})$, and
\item[(II)] $Y_n\setminus K(g_n)$ is disjoint from the level-$n$ bushes.
\end{itemize}
\emph{Status:} proved in the paper; the source traces the idea to Sullivan (\texttt{Feigen.tex}, lines 1679--1703).
\end{proposition}

\begin{theorem}[Primitive pull-off]
\label{thm:dl-fei-primitive}
For bounded $\bar p$ and small $\delta$, and for $n$ sufficiently large, if $\operatorname{Width}_\bullet(F_{n,i})=K\gg1$,
then either $\operatorname{Width}_\bullet(F)\ge 2K$, or the $n$-th renormalization is satellite and $F_{n-1}$ has a
$\delta$-almost-periodic rectangle of width at least $K/20$ between neighbouring satellite bushes.  \emph{Status:}
proved in the paper, refining Kahn I (\texttt{Feigen.tex}, lines 2640--2665).
\end{theorem}

\begin{theorem}[Satellite pull-off]
\label{thm:dl-fei-satellite}
For fixed $\bar p$ and sufficiently small $\delta$ there is $C_\delta\to\infty$ as $\delta\to0$ such that a sufficiently
wide $\delta$-almost-periodic rectangle implies $\operatorname{Width}_\bullet(F_2)\ge C_\delta K$ or
$\operatorname{Width}_\bullet(F)\ge C_\delta K$, with the analogous shifted alternative for $F=\mathcal R^{n\bullet}(G)$.
\emph{Status:} proved in the paper (\texttt{Feigen.tex}, lines 2877--2908).
\end{theorem}

\begin{theorem}[$\psi^\bullet$-bounds]
\label{thm:dl-fei-psi}
For every $\bar p$ there are $n>1$ and $K_{\bar p}>1$ such that
$\operatorname{Width}_\bullet\bigl[(\mathcal R^{n\bullet})^m(F)\bigr]\le K_{\bar p}$ for all sufficiently large $m$.
\emph{Status:} proved in the paper, using both pull-off theorems and Proposition~\ref{prop:dl-fei-teich}
(\texttt{Feigen.tex}, lines 3214--3221).
\end{theorem}

\begin{theorem}[Return to ordinary quadratic-like maps]
\label{thm:dl-fei-ql}
For every $\bar p$ there is $K_{\bar p}>1$ such that, for an infinitely renormalizable quadratic-like map $f$ of bounded
type $\bar p$ and $n\gg1$, the map $f_n$ has a quadratic-like restriction $X_n\to Y_n$ with
$\operatorname{Width}(Y_n\setminus X_n)\le K_{\bar p}$ (\texttt{Feigen.tex}, lines 3253--3254).  The transfer from
$\psi^\bullet$- to ordinary quadratic-like maps uses lemmas stated at lines 3260--3271.
\end{theorem}

\paragraph{Proof strategy (from the source).}
Primitive pull-off comes from thin--thick decompositions, Hubbard-tree alignment, the Grötzsch inequality and the
Covering Lemma (\texttt{Feigen.tex}, lines 430--460).  Satellite periodic rectangles are treated separately, with
spiralling numbers and wave estimates (lines 2908--2970).  Teichm\"uller contraction controls the growth of
prerenormalization widths (lines 1679--1703).  The two pull-off alternatives, with $C_\delta\gg\Delta_{\bar p}$, give
the uniform $\psi^\bullet$-bounds (lines 3214--3248).  The final statements use the Rigidity Theorem for MLC and
universality for the renormalization conclusions (lines 330--360).

\paragraph{External inputs.}
Lyubich, \emph{Dynamics of quadratic polynomials I--II} (rigidity for Main B; \texttt{Feigen.tex}, lines 3370--3374);
Lyubich, \emph{Feigenbaum--Coullet--Tresser universality and Milnor's hairiness conjecture} (Theorem C; lines 3376--3380);
Sullivan, \emph{Quasiconformal homeomorphisms in dynamics, topology and geometry}, Proc.\ ICM (source of the
Teichm\"uller contraction; lines 3415--3418); Kahn I (lines 3362--3365); Kahn--Lyubich, \emph{The quasi-additivity
law in conformal geometry}, Ann.\ of Math.\ 169 (2009) (Covering Lemma; lines 3367--3370); the satellite paper
(lines 3340--3344).

\paragraph{MLC scope.}
The claim is MLC at every Feigenbaum parameter in the paper's sense, that is, at every infinitely renormalizable
quadratic polynomial of bounded type, with local connectivity of the corresponding Julia sets
(\texttt{Feigen.tex}, lines 150--168 and 330--335).  The period-doubling point and the complex tripling renormalizations
are the most prominent cases (abstract).  The claim does not cover unbounded combinatorics, and it is not a proof of MLC at all
parameters.

\section{Kahn--Kapiamba--Lyubich: parabolic families}
\label{sec:kkl}

\paragraph{Identity.}
J.~Kahn, A.~Kapiamba and M.~Lyubich, \emph{MLC for parabolically bounded primitive renormalization}, arXiv:2606.27272
(2026) \cite{KahnKapiambaLyubich}.  The main file is \texttt{refs/arxiv-2606.27272/waves-with-Misha.tex} (metadata in
\texttt{00README.json}, lines 6--14).

\paragraph{Definitions.}
\begin{itemize}
\item \emph{Types and prime types.}  A type is a small Mandelbrot set $M\subsetneq\M$; it is \emph{prime} if no type
$M'$ satisfies $M\subsetneq M'\subsetneq\M$ (\texttt{waves-with-Misha.tex}, lines 111--115).  Prime types attached to
the main hyperbolic component are satellite; all others are primitive (lines 126--132).
\item \emph{Parabolic tail.}  $\{\Phi_n\circ\dots\circ\Phi_q(M):\ n\ge q\}$, where $M\subset\widehat M_{1/q}$ and the
$\Phi_j$ are the selected limb insertions (lines 144--151).
\item \emph{Parabolically bounded family.}  A finite union of types and parabolic tails (lines 149--153).  Equivalently,
after deleting finitely many types, its renormalizations have rotation number $1/q$, each limb contains at most $B$
small filled Julia components, and components in $\mathcal L_0$ map into $\bigcup_{|j|\le B}\mathcal L_j$
(lines 659--668).
\item \emph{Beau bounds} as in Definition~\ref{def:apriori-beau} (lines 117--123).
\item \emph{Rigid family.}  MLC holds at every infinitely renormalizable quadratic polynomial in $\M$ whose combinatorics
lie in the family (lines 160--163).
\item \emph{Super-optimism.}  For every $M>0$ there is a finite $\mathcal C'\subset\mathcal C$ such that every renormalizable
$\psi$-ql map $F$ with combinatorics in $\mathcal C\setminus\mathcal C'$ and $\mathcal W(\mathcal RF)\ge M$ satisfies
$\mathcal W(F)\ge M\,\mathcal W(\mathcal RF)$ (lines 207--215 and 611--612).
\end{itemize}

\paragraph{Results.}

\begin{theorem}[Theorem A]
\label{thm:kkl-A}
Every parabolically bounded family of prime primitive types has beau bounds.  \emph{Status:} proved in the paper
(\texttt{waves-with-Misha.tex}, lines 153--157).  The paper states that beau bounds alone do not establish MLC.
\end{theorem}

\begin{theorem}[Theorem B]
\label{thm:kkl-B}
Every family of prime primitive types with beau bounds and bounded puzzle Teichm\"uller diameter is rigid
(lines 160--165).  \emph{Status:} the paper observes that the argument of Lyubich's 1997 paper already gives rigidity
under this weaker hypothesis.  We record it as derived from the rigidity input of Section~\ref{sec:apriori-framework}.
\end{theorem}

\begin{theorem}[Theorem C]
\label{thm:kkl-C}
Every parabolically bounded family of prime primitive types is rigid (lines 167--171).  \emph{Status:} conditional on
Theorem B and on an input from the unpublished manuscript \emph{Elephant puzzles} (Dudko--Kapiamba--Lyubich--Peterson,
in preparation), which gives bounded puzzle Teichm\"uller diameter for parabolic tails (lines 167--171).
\end{theorem}

\begin{theorem}[Theorem D]
\label{thm:kkl-D}
Every parabolically bounded family of prime primitive types is super-optimistic (lines 218--221 and 693--696).
\emph{Status:} proved in the paper.
\end{theorem}

\begin{theorem}[Bounded optimism, cited from Kahn I]
\label{thm:kkl-bounded-optimism}
There is a universal $\varepsilon>0$ such that every bounded family of primitive types is $\varepsilon$-optimistic
(lines 607--610).  \emph{Status:} cited from Kahn I (\texttt{K06} in the source).
\end{theorem}

\begin{theorem}[Asymptotic optimism]
\label{thm:kkl-asymptotic}
Every super-optimistic family of prime primitive types has beau bounds (lines 612--615).  \emph{Status:} proved in the
paper, using bounded optimism.
\end{theorem}

\begin{theorem}[Value theorem]
\label{thm:kkl-value}
If $F$ is primitively renormalizable and periodically repelling, with core eigenvalue $\lambda_F$, then there is a
$\lambda_F$-expanding valuation for $F$ (\texttt{value1.tex}, lines 4--19).  \emph{Status:} proved in the paper.
\end{theorem}

\begin{theorem}[Thick--thin decomposition]
\label{thm:kkl-thick-thin}
For a Riemann surface $\mathcal V$ with ideal marking $\mathcal I$,
\[
\mathcal W(\mathcal V,\mathcal I)-O\bigl(|\chi(\mathcal V)|+|\mathcal I|\bigr)\ \le\ 2\,\|\mathcal W_{\rm can}(\mathcal V,\mathcal I)\|\ \le\ \mathcal W(\mathcal V,\mathcal I).
\]
\emph{Status:} cited from \texttt{KL:eyes}, which is Part IV, Section~\ref{sec:kahn-IV} (\texttt{waves-with-Misha.tex}, lines
1964--1978).  The disk case is Theorem~\ref{thm:kahn-IV-main-estimate}.
\end{theorem}

\begin{theorem}[Wave Lemma]
\label{thm:kkl-wave}
If $A'$ is a wave ending at $I'$ and $f$ has degree at most $d$ on $I'$, then some wave $A\subset f_*A'$ satisfies
$\mathcal W(A')\le\mathcal W(A)+O(1+\log d)$ (\texttt{waves-with-Misha.tex}, lines 2148--2157).  \emph{Source of the
notion:} the wave was introduced in the Dudko--Lyubich Feigenbaum paper, and the Wave Lemma is stated there
(\texttt{waves-with-Misha.tex}, lines 262--264).  The lemma is therefore cited from that paper.
\end{theorem}

\begin{theorem}[Wanderers theorem]
\label{thm:kkl-wanderers}
Let $g\colon X\to Y$ be proper holomorphic of degree $D$, let $\mathcal B\subset Y$ be a finite disjoint compact connected
set containing the critical values, and let $\mathcal A$ run over components of $g^{-1}(\bigcup\mathcal B)$.  Given an
oriented partial foliation $\mathcal F$ whose leaves start in $\widehat{\mathcal A}\subset\mathcal A$, end in $\mathcal A$,
visit at least $n-1$ intermediate components and do not cross their starting component, with degree at most $d$ over
$\mathcal B$,
\[
\|\mathcal W(\mathcal F)\|\ \le\ \frac{2d}{n}\,\|\mathcal W^{h+v}_{\rm can}(Y\setminus\textstyle\bigcup\mathcal B)\|
\;+\;C\bigl(|\chi(Y)|,|\mathcal B|,D\bigr).
\]
\emph{Status:} proved in the paper (\texttt{waves-with-Misha.tex}, lines 2342--2361).
\end{theorem}

\paragraph{Proof strategy (from the source).}
Parabolic tails are converted into uniformly parabolically bounded combinatorics after deleting finitely many types
(lines 670--684).  A life-improvement theorem says that high-type near-degeneracy forces width amplification
(lines 699--703).  Canonical laminations are analysed with the thick--thin decomposition, values, waves and translation
geometry (lines 224--267).  Theorem D follows from life improvement and parabolic boundedness (lines 704--710), Theorem A
from Theorem D and bounded optimism (lines 212--220), and Theorem C from A, B and the bounded puzzle Teichm\"uller diameter
of parabolic tails (lines 163--171).

\paragraph{External inputs.}
Kahn I (bounded optimism and the $\psi$-ql machinery); Kahn--Lyubich II (the asymptotic-optimism precedent; lines 626--630);
Kahn--Lyubich, \emph{The quasi-additivity law and the covering lemma} (covering lemma; lines 224--225 and 2342--2345);
Lyubich, \emph{Dynamics of quadratic polynomials I--II} (rigidity; puzzle renormalization; lines 160--163 and 344--347);
Dudko--Lyubich, Feigenbaum paper (waves and the Wave Lemma; lines 262--264); Kahn--Lyubich IV (thick--thin; \texttt{KL:eyes});
\emph{Elephant puzzles} (unpublished; used only in Theorem C).  The source also remarks that combining its results with the
Feigenbaum paper would give beau bounds and rigidity for families with finitely many primitive or satellite types, but this
needs a transfer between quadratic-like and $\psi^\bullet$-quadratic-like maps that is outside its scope (lines 174--175).

\paragraph{MLC scope.}
The claim is MLC at every infinitely renormalizable quadratic polynomial in $\M$ whose combinatorics belong to a parabolically
bounded family of prime primitive types, equivalently that these families are rigid (lines 160--171).  Theorem A alone does
not establish MLC (lines 160--162).  Theorem C depends on the unpublished input.

\section{Limbs and homeomorphisms (Dudko--Schleicher)}
\label{sec:dudko-schleicher}

\paragraph{Identity.}
D.~Dudko and D.~Schleicher, \emph{Homeomorphisms between limbs of the Mandelbrot set}, Proc.\ Amer.\ Math.\ Soc.\ 140
(2012), no.~6, 1947--1956 (electronically published 23 September 2011) \cite{DudkoSchleickerLimbs}.  The file is
\texttt{refs/Dudko\_S0002-9939-2011-11047-5.pdf}.

\paragraph{Context.}
The introduction recalls that an earlier result of the paper's reference $[BF_1]$ shows that limbs with equal denominators
are homeomorphic, and that those homeomorphisms extend to neighbourhoods of the limbs (\texttt{limbs-2011.txt}, Introduction).
The main theorem below preserves periods of hyperbolic components, and for that reason the source notes that it cannot extend
to neighbourhoods of the limbs.

\begin{theorem}[Main Theorem]
\label{thm:ds-main}
For any hyperbolic component of $\M$, let $\mathcal L_{p/q}$ and $\mathcal L_{p'/q}$ be two limbs with equal denominators.  Then
there exists a homeomorphism between them that preserves periods of hyperbolic components.  \emph{Status:} proved in the paper
(\texttt{limbs-2011.txt}, lines 50--54).  The paper relies on the decoration theorem and on Schleicher's combinatorial
results (Theorem~\ref{thm:ds-admissibility} below).
\end{theorem}

\begin{theorem}[Admissibility of numerators, Schleicher]
\label{thm:ds-admissibility}
If an angled internal address describes a combinatorial class in $\M$, then the numerators $p_k$ can be changed arbitrarily
(coprime to $q_k$) and the modified angled internal address still describes a combinatorial class in $\M$.  \emph{Status:} cited,
Schleicher's Theorem~2.3 as stated in the paper (\texttt{limbs-2011.txt}, lines 292--299).
\end{theorem}

\paragraph{The combinatorial model.}
There is a combinatorial model $\M_{\rm comb}$ of the Mandelbrot set, described as a limit of the arrangements $\M_n$ of hyperbolic
components of period up to $n$ in a certain sense, as in Douady's construction cited by the paper.  There is a canonical continuous
projection $\pi\colon\M\to\M_{\rm comb}$ whose fibres $\pi^{-1}(c)$ are compact, connected and full (\texttt{limbs-2011.txt}, lines 56--62).
The MLC conjecture ``can be stated as saying that $\pi$ is a homeomorphism'' (lines 62--63).  So the paper's framework expresses
MLC as the statement that $\pi$ is a homeomorphism.

\begin{theorem}[Canonical lifting]
\label{thm:ds-lifting}
Let $L_1,L_2\subset\M$ be two closed connected combinatorial sets.  For every continuous map
$f'\colon\pi(L_1)\to\pi(L_2)$ that is canonical with respect to small copies of $\M$, there is a continuous map
$f\colon L_1\to L_2$ such that the diagram with $\pi$ commutes (\texttt{limbs-2011.txt}, lines 447--450).  \emph{Status:} stated in
the paper; canonicity of $f'$ is defined there with respect to small copies (lines 431--437).
\end{theorem}

\paragraph{MLC scope.}
The paper proves a homeomorphism statement between limbs and a canonical lifting statement between combinatorial sets.  It does
not claim that limbs are locally connected, and it does not prove MLC: MLC is equivalent to $\pi$ being a homeomorphism.

\section{Summary of this chapter}
\label{sec:apriori-summary}

\begin{table}[ht]
\centering\small
\begin{tabularx}{\textwidth}{@{}>{\raggedright\arraybackslash}p{0.2\textwidth}>{\raggedright\arraybackslash}X>{\raggedright\arraybackslash}p{0.27\textwidth}@{}}
\textbf{Source} & \textbf{Result (label in this chapter)} & \textbf{Status in the source}\\
\hline
Kahn I &
Length comparison; bounded primitive type $\Rightarrow$ a priori bounds (Theorems~\ref{thm:kahn-I-length}, \ref{thm:kahn-I-main}) &
proved\\
Kahn II &
Decoration condition $\Rightarrow$ a priori bounds; beau bounds in decoration classes (Theorems~\ref{thm:kahn-II-main}, \ref{thm:kahn-II-beau}) &
proved, using the cited bounded-period theorem of Kahn I\\
Kahn III &
Molecule condition $\Rightarrow$ a priori bounds; beau bounds for molecule classes (Theorems~\ref{thm:kahn-III-main}, \ref{thm:kahn-III-beau}) &
proved, using the cited bounded-period theorem of Kahn I\\
Kahn IV &
Main estimate, bounded degree, amplification (Theorems~\ref{thm:kahn-IV-main-estimate}--\ref{thm:kahn-IV-amplification}) &
proved; no numbered main theorem printed in the TeX source\\
Dudko--Lyubich, satellites &
$\M$ and $J_{c_n}$ locally connected at constructed bounded-satellite parameters; $J_{c_n}$ of positive area (Theorem~\ref{thm:dl-sat-main}) &
proved; positive area uses the cited Avila--Lyubich estimates\\
Dudko--Lyubich, Feigenbaum &
Main A: beau bounds at Feigenbaum parameters; Main B: MLC at all of them (Theorems~\ref{thm:dl-fei-A}, \ref{thm:dl-fei-B}) &
Main A proved; Main B deduced from Main A and Lyubich's rigidity\\
Dudko--Lyubich, Feigenbaum &
Theorem C: renormalization conjecture for each $\bar p$ (Theorem~\ref{thm:dl-fei-C}) &
derived from cited inputs\\
KKL &
Theorems A and D: beau bounds and super-optimism for parabolically bounded families (Theorems~\ref{thm:kkl-A}, \ref{thm:kkl-D}) &
proved\\
KKL &
Theorems B and C: rigidity under beau bounds and bounded puzzle Teichm\"uller diameter (Theorems~\ref{thm:kkl-B}, \ref{thm:kkl-C}) &
B derived from Lyubich's argument; C conditional on unpublished \emph{Elephant puzzles}\\
Dudko--Schleicher &
Homeomorphisms between limbs with equal denominators, preserving periods (Theorem~\ref{thm:ds-main}) &
proved\\
\end{tabularx}
\caption{Results digested in this chapter and their status in the sources.}
\label{tab:apriori-summary}
\end{table}

\paragraph{Logical structure.}
The a priori bounds results form a chain.  Kahn I supplies the bounded-period improvement and the optimism machinery used by
Kahn II, Kahn III and KKL.  Dudko--Lyubich's Feigenbaum paper supplies the waves and the $\psi^\bullet$-bounds used by KKL.  The
thick--thin decomposition of Part IV is used by KKL.  The a priori bounds sources then pass to local connectivity through the
rigidity input of Section~\ref{sec:apriori-framework}.  The satellite paper proves local connectivity at its constructed
parameters by a different route, through pacman renormalization.  The limbs paper is independent of these arguments.

\paragraph{Connection to the Lean development.}
None of the theorems digested in this chapter is formalized.  The Lean development, described in
Chapter~\ref{chap:lean-inventory}, covers Green functions, puzzles, renormalization interfaces, the categorical branch,
certified approximations, and the elementary quadratic results (symmetries, first components, real slice, examples,
and multiplier maps).  The root statement is still conditional on the explicit \texttt{RootInput} hypothesis.

\paragraph{Not covered here.}
Hu--Jiang and Jiang are cited but not digested; their sources are not on arXiv.  Inou--Shishikura and Gaidashev--Yampolsky are cited
elsewhere and are not digested.  The unpublished \emph{Taming the elephants} and \emph{Elephant puzzles} are not available.  Lyubich's
Astérisque 261 paper (\emph{Dynamics of quadratic polynomials III}, bibliography key \texttt{LyubichIII}) is in the bibliography but is
digested in Chapter~\ref{ch:digests-covering-parapuzzle}.
